\documentclass[12pt, letterpaper]{article}
\usepackage{listings}
\usepackage{xcolor}
\usepackage{hyperref}
\usepackage{graphicx}

\lstset{
    basicstyle=\ttfamily\small,
    breaklines=true,
    frame=single,
    numbers=left,
    numberstyle=\tiny,
    keywordstyle=\color{blue},
    commentstyle=\color{gray},
    stringstyle=\color{red}
}

\title{Practical Report - File Transfer over TCP/IP in CLI}
\author{Nguyen Pham Truong An - 23BI14004}
\date{\today}

\begin{document}

\maketitle

\section{Protocol Design}

\subsection{Message Protocol Framework}
The system implements a framed protocol to distinguish between chat messages and file transfers, ensuring stream integrity and preventing data corruption.

\begin{verbatim}
+---------------------------------------------------+
|              PROTOCOL STRUCTURE                   |
+---------------------------------------------------+
| Message Protocol:                                 |
| +----------+-------------------------+            |
| | "MSG:"   | Message Content         |            |
| +----------+-------------------------+            |
|                                                   |
| File Transfer Protocol:                           |
| +----------+----------+----------+-------------+  |
| | "FILE:"  | filename | filesize | file data   |  |
| +----------+----------+----------+-------------+  |
|            (UTF-8)    (long)     (binary)        |
|                                                   |
| Acknowledgment Protocol:                          |
| +----------+                                      |
| | "ACK"    |                                      |
| +----------+                                      |
+---------------------------------------------------+
\end{verbatim}

\textbf{Protocol Types:}
\begin{itemize}
    \item \texttt{MSG:} - Text message with content prefixed by \texttt{MSG:}
    \item \texttt{FILE:} - File transfer initiation marker
    \item \texttt{ACK} - Acknowledgment for successful file receipt
    \item \texttt{/quit} - Graceful disconnection command
    \item \texttt{/send <filepath>} - User command to initiate file transfer
\end{itemize}

\section{System Architecture}

\subsection{Client-Server Organization}
The system implements a bidirectional, one-to-one TCP connection with multi-threaded communication.

\begin{verbatim}
+---------------------------------------+
|              SERVER                   |
+---------------------------------------+
| ServerSocket (bind, listen)           |
|        |                              |
|        v                              |
|   accept() -> Client Socket           |
|        |                              |
|        +-------------+                |
|        |             |                |
|   Send Thread    Receive Thread       |
|   (User Input)   (Network Listen)     |
|        |             |                |
|        v             v                |
|   DataOutputStream  DataInputStream   |
+---------------------------------------+

+---------------------------------------+
|              CLIENT                   |
+---------------------------------------+
| Socket (connect to server)            |
|        |                              |
|        v                              |
|   Established Connection              |
|        |                              |
|        +-------------+                |
|        |             |                |
|   Send Thread    Receive Thread       |
|   (User Input)   (Network Listen)     |
|        |             |                |
|        v             v                |
|   DataOutputStream  DataInputStream   |
+---------------------------------------+
\end{verbatim}

\textbf{Key Components:}
\begin{itemize}
    \item \textbf{Server}: Listens on port 8080, accepts one client connection, handles bidirectional communication
    \item \textbf{Client}: Connects to server via hostname resolution (\texttt{gethostbyname}), establishes TCP socket
    \item \textbf{Threading Model}: Each endpoint runs two threads - one for sending (user input) and one for receiving (network data)
    \item \textbf{Socket Operations}: Implements all core operations - socket(), bind(), listen(), accept(), connect(), send(), recv()
\end{itemize}

\section{File Transfer Implementation}

\subsection{Transfer Flow with Atomic Operations}

The file transfer implements a three-phase protocol ensuring atomicity:

\begin{verbatim}
PHASE 1: TRANSMISSION
  Sender                    Receiver
    |                          |
    |---- "FILE:" ------------>|
    |---- filename ----------->| (UTF-8 string)
    |---- filesize ----------->| (long, 8 bytes)
    |---- file data ---------->| (binary stream)
    |                          |
    |                          | Write to "recv_filename"
    
PHASE 2: ACKNOWLEDGMENT
    |                          |
    |<------- "ACK" -----------| File successfully received
    
PHASE 3: FINALIZATION
    | Delete original file     |
    |                          | Rename: recv_filename -> filename
\end{verbatim}

\subsection{Core File Transfer Code}

\begin{lstlisting}[language=Java,caption=Client/Server File Sending,label=code:sendfile]
private void sendFile(String path) {
    try {
        File file = new File(path);
        if (!file.exists()) {
            System.out.println("File not found: " + path);
            return;
        }
        
        // Send file transfer protocol header
        dos.writeUTF("FILE:");
        dos.writeUTF(file.getName());
        dos.writeLong(file.length());
        
        // Stream file data in 4KB chunks
        FileInputStream fis = new FileInputStream(file);
        byte[] buffer = new byte[BUFFER_SIZE];
        int read;
        while ((read = fis.read(buffer)) != -1) {
            dos.write(buffer, 0, read);
        }
        dos.flush();
        fis.close();
        
        // Wait for acknowledgment
        String ack = dis.readUTF();
        if (ack.equals("ACK")) {
            if (file.delete()) {
                System.out.println("Transferred: " + file.getName());
            } else {
                System.out.println("Sent but failed to delete: " 
                    + file.getName());
            }
        }
    } catch (IOException e) {
        System.err.println("Send failed: " + e.getMessage());
    }
}
\end{lstlisting}

\begin{lstlisting}[language=Java,caption=Client/Server File Receiving,label=code:receivefile]
private void receiveFile() {
    try {
        // Read file metadata
        String filename = dis.readUTF();
        long size = dis.readLong();
        
        // Write to temporary file with recv_ prefix
        FileOutputStream fos = new FileOutputStream("recv_" + filename);
        byte[] buffer = new byte[BUFFER_SIZE];
        long remaining = size;
        int read;
        
        // Read exact file size to prevent over-reading
        while (remaining > 0 && 
               (read = dis.read(buffer, 0, 
                   (int)Math.min(buffer.length, remaining))) != -1) {
            fos.write(buffer, 0, read);
            remaining -= read;
        }
        fos.close();
        
        // Send acknowledgment
        dos.writeUTF("ACK");
        dos.flush();
        
        // Atomic rename: recv_filename -> filename
        File tempFile = new File("recv_" + filename);
        File finalFile = new File(filename);
        if (tempFile.renameTo(finalFile)) {
            System.out.println("Received: " + filename);
        } else {
            System.out.println("Received but rename failed: recv_" 
                + filename);
        }
    } catch (IOException e) {
        System.err.println("Receive failed: " + e.getMessage());
    }
}
\end{lstlisting}

\subsection{Socket Operations Mapping}

\begin{table}[h]
\centering
\begin{tabular}{|l|l|}
\hline
\textbf{Socket Operation} & \textbf{Java Implementation} \\ \hline
\texttt{socket()} & \texttt{new Socket()}, \texttt{new ServerSocket()} \\ \hline
\texttt{setsockopt()} & \texttt{setReuseAddress()}, \texttt{setSoTimeout()} \\ \hline
\texttt{bind()} & \texttt{serverSocket.bind(InetSocketAddress)} \\ \hline
\texttt{gethostbyname()} & \texttt{InetAddress.getByName(host)} \\ \hline
\texttt{listen()} & Implicit in \texttt{ServerSocket} creation \\ \hline
\texttt{connect()} & \texttt{socket.connect(InetSocketAddress)} \\ \hline
\texttt{accept()} & \texttt{serverSocket.accept()} \\ \hline
\texttt{send()} & \texttt{DataOutputStream.write()}, \texttt{writeUTF()} \\ \hline
\texttt{recv()} & \texttt{DataInputStream.read()}, \texttt{readUTF()} \\ \hline
\end{tabular}
\caption{Socket API to Java Mapping}
\end{table}

\section{Key Implementation Features}

\subsection{Atomic File Transfer}
\begin{itemize}
    \item \textbf{Temporary Prefix}: Files written with \texttt{recv\_} prefix during transfer
    \item \textbf{Acknowledgment Protocol}: Sender only deletes after receiving \texttt{ACK}
    \item \textbf{Atomic Rename}: Receiver renames file only after successful write and ACK sent
    \item \textbf{Safety}: Prevents partial files and ensures transaction-like behavior
\end{itemize}

\subsection{Protocol Framing}
\begin{itemize}
    \item \textbf{Clear Separation}: \texttt{MSG:} prefix distinguishes text from binary data
    \item \textbf{Size Tracking}: File size header prevents stream corruption
    \item \textbf{UTF-8 Encoding}: Filename transmitted as UTF-8 for international support
    \item \textbf{Binary Integrity}: Raw byte transfer after headers maintains file integrity
\end{itemize}

\subsection{Error Handling}
\begin{itemize}
    \item \textbf{File Validation}: Checks file existence before transmission
    \item \textbf{Connection Resilience}: Handles network failures gracefully
    \item \textbf{Resource Management}: Proper stream closure in finally blocks
    \item \textbf{Timeout Configuration}: 30-second socket timeout prevents hanging
\end{itemize}

\subsection{Concurrency Model}
\begin{itemize}
    \item \textbf{Dual Threading}: Separate threads for sending and receiving
    \item \textbf{Volatile Flag}: \texttt{volatile boolean running} for thread-safe shutdown
    \item \textbf{Non-blocking Input}: User can type while receiving messages/files
    \item \textbf{Thread Coordination}: Clean shutdown via \texttt{/quit} command
\end{itemize}

\section{Usage Examples}

\subsection{Compilation and Execution}
\begin{verbatim}
# Compile
javac Server.java Client.java

# Run Server (Terminal 1)
java Server 8080

# Run Client (Terminal 2)
java Client 127.0.0.1 8080
\end{verbatim}

\subsection{Sample Session}
\begin{verbatim}
Server Terminal:
$ java Server 8080
Server listening on port 8080
Client connected: 127.0.0.1
Hello from server!
Client: Hi server!
/send document.pdf
Transferred: document.pdf
Received: report.txt

Client Terminal:
$ java Client 127.0.0.1 8080
Connecting to 127.0.0.1:8080
Connected
Server: Hello from server!
Hi server!
Received: document.pdf
/send report.txt
Transferred: report.txt
\end{verbatim}

\section{Conclusion}

This implementation demonstrates a complete TCP/IP file transfer system with the following achievements:

\begin{itemize}
    \item Full implementation of core socket operations
    \item Atomic file transfer with acknowledgment protocol
    \item Bidirectional communication with chat capability
    \item Thread-safe concurrent send/receive operations
    \item Robust error handling and resource management
    \item Clean protocol framing preventing stream corruption
\end{itemize}

The system successfully demonstrates TCP/IP concepts while providing practical file transfer functionality in a command-line environment.

\end{document}